% This file recreates the opening four-page teaching sequence. It introduces the
% physical components, walks through one diplomacy quest and one combat quest, and
% uses labeled placeholders wherever the source depends on product photography.

% The source cover occupies the upper portion of page one rather than a separate
% leaf, so the title treatment and first rules section intentionally share a page.
\begin{center}
  {\sffamily\bfseries Questweight Studios}\par
  {\scriptsize By: Jeremy Offutt \& Curtis Mortensen}\par
  \vspace{4pt}
  \DndrGraphicPlaceholder[0.62\linewidth]{1.45in}
  \vspace{-3pt}
  {\Huge\bfseries\color{DndrGold} Dragons \& Dice Rolls}\par
  {\large\bfseries Heroic Chaos Without the Homework}\par
  {\small d4 \quad d6 \quad d8 \quad d10 \quad d12 \quad D20}
\end{center}

\begin{multicols}{2}
\DndrSectionBanner{Section 1: Quickstart Guide}

\subsection*{1.1 Welcome to Beginner's School}
Dice \& Dragons is a fast-paced, cooperative storytelling game where you and your friends compose an epic adventure together. It can be as ridiculous or serious as your table wants it to be. If you have ever wanted the best parts of traditional tabletop roleplaying games - slaying dragons, delivering dramatic speeches, and collecting cool loot - without committing years of your life or performing endless homework, you are in the right place.

\DndrCallout{Golden Rule of Tone}{Choose a shared tone: comedy, epic fantasy, or anything in between. Not as a rule - just a way to keep everyone aiming for the same vibe.}

\subsection*{1.2 Automated Setup for Your First Adventure}
To learn the game by playing immediately, set it up as in the picture. Put all decks of cards in the center of the table.
\DndrGraphicPlaceholder{0.72in}

Pass the Character Mats and Backgrounds around the table. Each player chooses exactly one of each:
\begin{enumerate}
  \DndrStep{Select a Class Mat:}{Choose from Force of Nature, Tactician, Rapscallion, Meat Shield, Pack Leader, or Minstrel. Grab the specific class dice assigned to your Class Mat line.}
  \DndrStep{Select a Background Card:}{Place it directly onto its designated slot on your Class Mat.}
  \DndrStep{Fill Your Hearts:}{Take the starting Heart tokens shown on your Class Mat and place them into the heart-shaped slots.}
  \DndrStep{Acquire Your Signature Item:}{Deal 3 Item cards to each player. Replace any red- or green-bordered card; only basic items can be signature items at the start. Choose one permanent, iconic Signature Item and discard the other two.}
\end{enumerate}

\DndrCallout{Choosing Strategy}{Choose the most versatile item - one useful in combat, diplomacy, and puzzles. Barring that, choose a combat item.}

\begin{enumerate}[start=5]
  \DndrStep{Draw Your Starting Gear:}{Draw one additional Item card. If its skill tags match a skill printed on your Class Mat or Background, keep it. Otherwise return it and draw until you find an item you can legally wield.}
\end{enumerate}

\DndrCallout{Backpack Comfort}{Equip up to 4 Items at a time. Your Signature Item does not count against this limit. Unequipped Items remain in your unlimited backpack.}

% A physical page break here mirrors the source: page one teaches setup, while
% page two starts with character introductions and the first complete Quest flow.
\end{multicols}
\clearpage
\begin{multicols}{2}

\subsection*{1.3 Introducing Your Character}
Pick a Name and draw a Flashback Prompt. Introduce yourself in character: How did an ordinary person with your Background become a legendary member of your Class, and how did you acquire your Signature Item? Use your Prompt for inspiration and tell the story in first person.

When another player finishes, applaud their effort. Give them one of your color's Applause Tokens. Helping them does not hurt you.

\subsection*{1.4 The Opening Scene \& Choosing a Goal}
Choose a Party Goal. As a group, decide what grand objective your adventure is trying to accomplish. This goal keeps everyone moving in the same direction.

The youngest player becomes the first Architect, serving as narrator and game master for the opening quest. Give them the Architect Token.

\subsection*{1.5 Quest 1: Secure a Happy Ever After}
\begin{enumerate}
  \DndrStep{Read the introduction:}{The Architect draws ``Secure a Happy Ever After'' and reads the description aloud.}
  \DndrStep{Read the first scenario:}{Help the royal lover sort out their complicated feelings to reveal the NPC the party must persuade.}
  \DndrStep{Set the Challenge:}{Draw 1-3 Challenge Cards and hand them to players. Each recipient explains how their Challenge makes the Quest harder.}
  \DndrStep{Deal:}{Give every player except the Architect 3 white Diplomacy Cards. The Architect draws and reads a blue NPC Card.}
  \DndrStep{Submit:}{Each player secretly chooses 1 Diplomacy Card. Add random cards equal to the Challenge Rating, then shuffle the pile.}
  \DndrStep{Activate an Item (optional):}{Spend 1 Applause Token, describe a brief flashback about the gear, and resolve its effect.}
  \DndrStep{The Architect's Verdict:}{Reveal and read every submitted card. Choose the outcome that best fits the negotiation.}
\end{enumerate}

\textbf{Success:} If a player's card is chosen, the party persuades the NPC. Narrate the success and how it leads to the next Quest.

\textbf{Failure:} If an added card is chosen, competing interests, misinformation, or bad luck win the day. Narrate what went wrong and how the party moves on.

\begin{enumerate}[start=8]
  \DndrStep{Rewards:}{On success, deal 1 Item Card to each player. Then draw Item Cards equal to twice the Challenge Rating and distribute them among the party.}
\end{enumerate}

% The source gives combat its own continuation page so its round procedure remains
% usable at the table instead of being squeezed beneath the Diplomacy instructions.
\end{multicols}
\clearpage
\begin{multicols}{2}

\subsection*{1.6 Running Combat: Fend Off Other Suitors}
\begin{enumerate}
  \DndrStep{Present the threat:}{Read ``Fend off other suitors.''}
  \DndrStep{Choose monsters:}{Draw or pick a monster to act as the suitors. Have fun explaining why.}
  \DndrStep{Set the Challenge:}{Draw 1-3 Challenge Cards. Players explain why those complications belong in the story.}
  \DndrStep{Build the encounter:}{Use the Monster Manual pages and the total difficulty on the Challenge Cards. Place the appropriate Hearts on each monster.}
  \DndrStep{Monsters attack first:}{Roll the attack die in the stat block and compare it to the target's Defense. Meeting the Defense deals the listed damage.}
  \DndrStep{Players act:}{In clockwise order, choose a Basic Attack, Activate an Item, or Charge.}
  \DndrStep{Continue rounds:}{After every player acts, the monsters attack again. Continue until all monsters or all heroes are defeated.}
  \DndrStep{Rewards:}{On success, use the same reward procedure as the Diplomacy Quest.}
\end{enumerate}

\DndrCallout{Heroic Actions}{Once per combat, when you would lose your final Heart, return every Applause Token you hold to its owner. Stay standing with half a Heart instead.}

\subsection*{1.7 Downtime and Rotating Roles}
After Quest 1, every player heals half a Heart, players may trade Items, and the Architect Token passes to the player on the left.

\subsection*{1.8 Part 2, Quest 1: Combat}
The new Architect leads the second part of the Quest. Monsters act first. On a player turn, a Basic Attack automatically deals half a Heart, an Item costs an Applause Token and a flashback, and Charge spends a token now to double the next attack's damage.

% Quest 2 begins on the fourth quickstart page, matching the photographic spread
% in the reference while leaving space for a future play-session image.
\end{multicols}
\clearpage
\begin{multicols}{2}

\subsection*{2.0 Quest 2: Help a Dragon Fall Asleep}
Pass the Architect Token. Run ``Help a Dragon Fall Asleep,'' featuring a Puzzle Challenge instead of diplomacy.
\DndrGraphicPlaceholder{2.35in}
\begin{enumerate}
  \DndrStep{Read the description:}{Perform a bedtime story.}
  \DndrStep{Draw Challenges:}{Pass out Challenge Cards and let players explain them.}
  \DndrStep{Set up the encounter:}{Play Knucklebones as described in Section 5.}
  \DndrStep{Narrate the ending:}{Explain how success advances the Party Goal or how failure creates a setback.}
  \DndrStep{Pass out rewards:}{Use the normal reward procedure.}
\end{enumerate}

\subsection*{2.1 Downtime}
Pass the Architect Token and trade Items if desired. Keep in mind that the next Quest is combat: despite your bedtime story, noisy creatures still need attention.

\subsection*{2.2 Combat}
Combat works the same as ``Fend Off Other Suitors.'' For edge cases, use Section 4.5.

\subsection*{2.3 Finish}
Congratulations! You have survived Beginner's School. You now know Quests, Flashbacks, Items, Challenges, Diplomacy, Puzzles, and Combat. Continue these heroes or start fresh with a new party.
\end{multicols}
